% SOURCE_AUTHORITY: sga5_fr_workpass.tex, lines 10035--10095 (LF-normalized unit).
% SOURCE_SHA256: 791F4EFFC5E02832D5D77ED03518C8156D6F07E4C8238B03545DB93D883FBB28
\emph{Demostración.} Como el morfismo $h$ es una composición de un número finito de morfismos del tipo
\[
P_T(F)\to T\qquad (F\text{ }\mathcal O_T\text{-módulo localmente libre}),
\tag{Q}
\]
de (2.2.1) se sigue que
\[
\mathrm Rh_*(A)\simeq\bigoplus_j R^jh_*(A)[-j].
\]
Por la misma razón, se tiene también
\[
\mathrm Rg_*(A)\simeq\bigoplus_j R^jg_*(A)[-j].
\]
Utilizando ([6] 2.16), se deduce un isomorfismo
\[
\mathrm Rf_*(g_*A)\simeq\bigoplus_j R^j f_*(g_*A)[-j].
\]
Como $g$ es liso con fibras geométricamente conexas, $A_X\simeq g_*(A_D)$, de modo que queda por demostrar que
\[
R^{2i+1}f_*(A)=0\qquad (i\in\mathbb N).
\tag{5.2.1}
\]
Para ello, por la conservatividad del sistema de funtores fibra (SGA 4 VIII 3.5), se reduce inmediatamente al caso en que $S$ es el espectro de un cuerpo separablemente cerrado $k$. Entonces, como el morfismo $g$ es una composición de fibrados proyectivos, de (2.2.4) se sigue que la imagen inversa
\[
H^{2i+1}(X,A_X)\to H^{2i+1}(D,A_D)
\]
es un monomorfismo, de modo que se puede suponer que $X=D$. En este caso, como el morfismo de proyección $D\to k$ es a su vez una composición de fibrados proyectivos, la afirmación se sigue por inducción de (2.2.4).

\begin{statement}{Corolario 5.3}
Para todo morfismo de esquemas $u:S\to T$, la sucesión espectral de Leray
\[
E_2^{pq}=R^p u_*\,R^q f_*(A)\Rightarrow R^{p+q}(u\circ f)_*(A)
\]
es degenerada.
\end{statement}

\emph{Demostración.} Cf. [6] 1.2.

Precisemos ahora la estructura multiplicativa de la cohomología de $X$, examinando primero un caso particular. Con las notaciones anteriores, se obtiene un morfismo de anillos
\[
\alpha:A^*(S)[T_{ij_i}]_{1\le i\le m,\,1\le j_i\le p_i}\to A^*(X)
\]
enviando $A^*(S)$ en $A^*(X)$ por la imagen inversa $f^*$ y, para todo par $(i,j)$ como antes, el elemento $T_{ij_i}$ en la $j_i$-ésima clase de Chern $c_{j_i}(E_i)$. Para todo entero $p$, se tiene, por (3.4 (iii)),
\[
c_p(f^*E)=\sum_{i_1+i_2+\cdots+i_m=p}c_{i_1}(E_1)\cdots c_{i_m}(E_m),
\]
de modo que, por la funtorialidad de las clases de Chern, el morfismo se anula sobre el ideal $J_X$ generado por los elementos
\[
c_p(E)-\sum_{i_1+\cdots+i_m=p}T_{1,i_1}\cdots T_{m,i_m}
\]
para $p$ variable. De aquí se obtiene un morfismo de anillos
\[
\varphi_X:\bigl(A^*(S)[T_{ij_i}]_{(i,j_i)}\bigr)/J_X\to A^*(X).
\]

\begin{statement}{Proposición 5.4}
El morfismo de anillos $\varphi_X$ anterior es un isomorfismo. El $A^*(S)$-módulo $A^*(X)$ es libre de rango
\[
\frac{r!}{p_1!p_2!\cdots p_m!}.
\]
\end{statement}
